\documentclass{scrartcl}
\usepackage{psetconfig}
\usepackage{bookmark}
\usepackage{scribe}

%%%%%%%%%%%%%%%%%%%%%%%%%%%%%%%%%%%%%%%%%%%%%
%Fill in the appropriate information below
\clearpairofpagestyles
\ihead{Songyu Ye}
\ohead{MATH 6110} 
\chead{\textbf{Problem Set 6}}
%%%%%%%%%%%%%%%%%%%%%%%%%%%%%%%%%%%%%%%%%%%%%

\begin{document}
\begin{problem}{3.24}
Suppose $F$ is an increasing function on $[a,b]$. \begin{enumerate}
    \item Prove that we can write \begin{align*}
        F = F_A + F_C + F_J
    \end{align*} where $F_A$ is absolutely continuous, $F_C$ is continuous but $F'_C(x)= 0$ for a.e. x, and $F_J$ is a jump function.
    \item Each component of the decomposition is unique up to an additive constant.
\end{enumerate}
\end{problem}

\begin{solution}
We can without loss of generality assume that $F$ is bounded on $[a,b]$.
Indeed, if $F$ is unbounded, we can consider the restriction of $F$ to $[a,b]\cap [-n,n]$ for $n\in\N$
and apply the following argument to each restriction. 
Then we can take the union of the decompositions of each restriction to obtain a decomposition of $F$.

Then knowing that $f$ is bounded and increasing, we can apply Lemma 3.13 to conclude the
existence of a decomposition $F = F_i + F_j$ where 
$F_i$ is increasing continuous and $F_j$ is a jump function.

Since $F_i$ is increasing it is of bounded variation, and hence by Corollary 3.7
it is locally integrable on $[a,b]$. Therefore 
we get another decomposition 
\begin{align*}
    F_i = F_C + F_A
\end{align*} where \begin{align*}
    F_A(x) &= \int_a^x F'(t)dt \\
    F_C(x) &= F_i - F_A
\end{align*} Then we can verify the following claims:
\begin{enumerate}
    \item $F_A$ is absolutely continuous. This is because $F_A$ is the integral of a locally integrable function.
    \item $F_C$ is continuous. This is because $F_C$ is the difference of two continuous functions.
    \item $F_C'(x) = 0$ for a.e. $x$. This follows from the Lebesgue differentiation theorem and the fact that $F_C$ is the difference of two increasing functions.
\end{enumerate}

Thus we have shown the existence of the decomposition $F = F_A + F_C + F_J$.
Let us now verify the uniqueness of the decomposition.

Suppose $F = F_A' + F_C' + F_J'$ is another decomposition. For notation, let $\Delta_A = F_A - F'_A$
and similarly for the other components $C,J$. Then we have \begin{align*}
    F_A' + F_C' + F_J' = F_A + F_C + F_J
\end{align*} and hence we have \begin{align*}
    \Delta_J = F_J - F_J' = F_A' + F_C' - F_A - F_C 
\end{align*} and the left hand side is a jump function, and the right hand side is continuous.
 Therefore we see that $\Delta_j = C_j$.

We also have\begin{align*}
    \Delta_A = F_A' - F_A = F_C - F_C' + F_J - F_J'
\end{align*} Since $F_A$ are absolutely continuous, the left hand side is. 
Then we can observe further that the left hand side is differentiable a.e. 
and thus an application of the Lebesgue integration theorem tells us that \begin{align*}
    \Delta_A(x) - \Delta_A(a) &= \int_a^x (F_C - F_C' + F_J - F_J')'(t)dt \\
    &= \int_a^x(F_A -F_A')'(t)dt = 0
\end{align*} and so we see that $\Delta_A$ is a constant function.

Finally we have \begin{align*}
    \Delta_C = \Delta_A + \Delta_J = C_A + C_J
\end{align*} and so we see that $\Delta_C$ is a constant function.
This concludes the proof of uniqueness.
\end{solution}

\begin{problem}{3.29}
    Let $\Gamma = \set{z(t)\st a\leq t \leq b}$ be a curve and suppose it satisfies Lipschitz with exponent $1/2\leq \alpha\leq 1$. Show that 
    $m(\Gamma^\delta) = O(\delta^{2-1/\alpha})$ for $0<\delta\leq 1$.
\end{problem}

\begin{solution}
This proof is not complete, but here is a sketch of how I am sort of considering attacking the problem. I also verify the claim in the case $\alpha = 1$ as a sanity check.


Let $\delta \in (0,1)$ arbitrary and reparametrize the curve $z:[9,1]\to\R^2$ and then partition the interval $0\leq t_0 < t_1 < \dots < t_n = 1$ so tha $t_{i+1} - t_i = c$. Then by the Lipschitz
condition we can see that \begin{align*}
    \abs{z(t_{i+1}- t_i)} \leq Ac^\alpha
\end{align*} Now consider balls of radius $\delta$ centered at each of the points $z(t_i)$. These balls cover a set $S\subset \R^2$ which is clearly contained in $\Gamma^\delta$.
Now I want to make some sort of statement that for a proper choice of $c$ which will certainly have to depend on $\delta$, $S = K_\delta$.

\hfill

Now if we did not know that $z$ was Lipschitz, this could not possibly be true unless we took a ball at every point along the curve. This is because for example to cover the $\Gamma^\delta$
corresponding to when $\Gamma$ is a straight line (i.e. this coincides with the case that $\alpha$ = 1) we need to take a ball at every point. In fact if we do not, then we can find a point $p$ which is not in the union of the balls but
which is distance $\delta$ away from the curve, namely the point distance $\delta$ away from the missing point along the curve in the direction of a normal vector to the curve.

\hfill 

With this Lipschitz condition we hope for a statement of the following nature. 

\textbf{Claim:} For an appropriate choice of $c$, then any $p\in\R^2$ which satisfies $d(p,z(t_i)) > \delta$ for all $t_i$ implies that $d(p,\Gamma) > \delta$.

If I am able to show this and actually get my hands on such a $c$ then one could proceed by actually computing the measure of the cover corresponding to such a choice
of $c$ and then see if it is in fact $O(\delta^{2-1/\alpha})$. This is where I sort of got stuck. 

\hfill

Here is a sort of sanity check that I carried out for the case $\alpha = 1$.

One can convince themself that $\alpha = 1$ as a global Lipschitz constant actually implies that $z$ is in fact linear and therefoore we can see that $\Gamma^\delta$ is in fact
just a rectangle with half circle ends, we can compute the measure of this and see that it is in fact $O(\delta) = O(\delta^{2-1/1})$.
\end{solution}

\begin{problem}{2}
    Suppose $I_1, I_2, \cdots$ are disjoint open intervals in $\R$. Then there are two finite subcollections 
    $\set{I_1', I_2', \cdots, I_{k_1}'}$ and $\set{I_1'', I_2'', \cdots, I_{k_2}''}$ such that \begin{align*}
        \bigcup_{i=1}^\infty I_i = \bigcup_{j=1}^{k_1} I_{1,j}' \cup \bigcup_{j=1}^{k_2} I_{2,j}''
    \end{align*} and each subcollection consists of disjoint intervals.

    Hint: Choose $I_1'$ to be an interval whose left end-point is as far left as possible. Discard all intervals contained in $I_1'$.
     If the remaining intervals are disjoint from $I_1'$, select again an interval as far to the left as possible, and call it  $I_2'$.
     Otherwise choose an interval that intersects $I_1'$ , but reaches out to the right as far as possible, and call this interval $I_2'$.
     Repeat this procedure.
\end{problem}

\begin{solution}
We proceed as the hint suggests, and choose $I_1'$ to be an interval whose left end-point is as far left as possible.
Discard all intervals contained in $I_1'$ and call the remaining intervals $I_1, I_2, \cdots,I_{k_1}$.
If the remaining intervals are disjoint from $I_1'$, select again an interval as far to the left as possible, and call it  $I_2'$ and repeat.
Otherwise choose an interval that intersects $I_1'$ , but reaches out to the right as far as possible, and call this interval $I_1''$.
Discard all intervals contained in $I_1''$ and repeat.

Now we can verify some properties of the construction. In particular we can observe that
\begin{enumerate}
    \item $I_1', I_2', \cdots$ are disjoint intervals. This is because we choose $I_1', I_2', \cdots$ to be intervals that are as far left as possible and we've removed all the ones in between.
    \item $I_1'', I_2'', \cdots$ are disjoint from $I_1'', I_2'', \cdots,I_{k_2}$. This is for the same reason as above.
    \item There are only fintely many $I_1', I_2', \cdots$ and $I_1'', I_2'', \cdots$. This is because we can only choose $I_1', I_2', \cdots$ and $I_1'', I_2'', \cdots$ finitely many times.
    \item The $I_1', I_2', \cdots$ and $I_1'', I_2'', \cdots$ cover, in particular \begin{align*}
        \bigcup_{i=1}^\infty I_i = \bigcup_{j=1}^{n_1} I_{1,j} \cup \bigcup_{j=1}^{n_2} I_{2,j}
    \end{align*}
\end{enumerate}
Therefore we have verified the desired properties of the construction.
\end{solution}

\begin{problem}{7}
    Consider the function 
    \begin{align*}
        f_1(x) = \sum_{n=0}^\infty 2^{-n}e^{2\pi\cdot i2^nx}
    \end{align*} \begin{enumerate}
        \item Prove that $f_1$ satisfies $\abs{f_1(x)-f_1(y)}\leq A_\alpha\abs{x-y}^\alpha$ for all $x,y\in\R$ and for each $\alpha\in(0,1)$.
        \item $f$ is nowhere differentiable, hence not of bounded variation.
    \end{enumerate}
\end{problem}

\begin{solution}
We compute \begin{align*}
    \abs{f_1(x) - f_1(y)} &\leq \sum_{n=0}^\infty 2^{-n}\abs{e^{2\pi\cdot i2^nx} - e^{2\pi\cdot i2^ny}} \\
    &= \sum_{n=0}^\infty 2^{-n}\abs{e^{2\pi\cdot i2^nx} - e^{2\pi\cdot i2^ny}}^{1-\alpha} \abs{e^{2\pi\cdot i2^nx} - e^{2\pi\cdot i2^ny}}^{\alpha} \\
    &\leq \sum_{n\geq 0}2^{-n}2^{1-\alpha}\abs{2\pi i 2^n}\abs{x-y}^\alpha \\
    &= 2^{1-\alpha}\abs{2\pi i}\abs{x-y}^\alpha \sum_{n\geq 0}2^{-\alpha n} \\
    &= A_\alpha\abs{x-y}^\alpha
\end{align*} where the second inequality is by the mean value theorem.
Specifically the mean value theorem applied to the function $f(t) = \exp(2\pi i2^nt)$ on the interval $[x,y]$ tells us that \begin{align*}
    \abs{e^{2\pi\cdot i2^nx} - e^{2\pi\cdot i2^ny}} &\leq \abs{2\pi\cdot i2^n(x-y)} \\
    &\leq \abs{2\pi\cdot i2^n(x-y)}
\end{align*}
Therefore we indeed see that $f_1$ satisfies $\abs{f_1(x)-f_1(y)}\leq A_\alpha\abs{x-y}^\alpha$ for all $x,y\in\R$ and for each $\alpha\in(0,1)$.

Now we show that $f_1$ is nowhere differentiable. 

Suppose $f_1$ is differentiable at $x_0$. Then we have \begin{align*}
    \lim_{x\to x_0}\frac{f_1(x)-f_1(x_0)}{x-x_0} &= \lim_{x\to x_0}\sum_{n=0}^\infty 2^{-n}\frac{e^{2\pi\cdot i2^nx} - e^{2\pi\cdot i2^nx_0}}{x-x_0} \\
    &= \sum_{n=0}^\infty 2^{-n}2\pi\cdot i2^nx_0  \\
    &= 2\pi\cdot i\sum_{n=0}^\infty 2^{-n}2^nx_0 \\
    &= 2\pi\cdot i\sum_{n=0}^\infty x_0 \\
    &= \infty
\end{align*} so $f_1$ is not differentiable at $x_0$. Since $x_0$ is arbitrary, $f_1$ is nowhere differentiable.
\end{solution}


\newpage

\begin{problem}{7.1}
Consider the linear map defined on functions of $d$ variables with values in functions on $d+1$ variables
given by \begin{align*}
    Tf(x,t) = \sum_{k,l\st k\leq 0, k+l\geq 0}2^{-(k+l)d}\langle f, 1_{[0,2^k]^d}\rangle 1_{[0,2^k]^d}(x)1_{[0,4^l]}(t)
\end{align*}
Prove that $T$ is a bounded linear operator from $L^p(\R^d)$ to $L^{p}(\R^{d+1})$ as long
as $p > 2(d+1)/d$.
\end{problem}
\begin{solution}
Let us see that $Tf$ is bounded linear operator. It is clear that $Tf$ is a linear operator. Let us see that $T$ is bounded.
The computations will also show that $Tf$ is $p$-integrable if $f$ is $p$ integrable. From the generalized Minkowski inequality we have \begin{align*}
    ||Tf||_p &\leq \sum_{k,l\st k\leq 0, k+l\geq 0}||2^{-(k+l)d}\langle f,1_{[0,2^k]^d}\rangle1_{[0,2^k]^d}(x)1_{[0,4^l]}(t)||\\
    &\leq \sum_{k,l\st k\leq 0, k+l\geq 0}|2^{-(k+l)d}\langle f,1_{[0,2^k]^d}\rangle|||1_{[0,2^k]^d}(x)1_{[0,4^l]}(t)|| \\
    &\leq\sum_{k,l\st k\leq 0, k+l\geq 0}2^{-(k+l)d}|\langle f,1_{[0,2^k]^d}\rangle|||1_{[0,2^k]^d}(x)1_{[0,4^l]}(t)|| 
\end{align*} Recall that $||fg||_1 \leq ||f||_p||f||_q$ for $1/p + 1/q = 1$. This implies that 
\begin{align*}
    ||Tf||_p &\leq \sum_{k,l\st k\leq 0, k+l\geq 0}2^{-(k+l)d}||f||_p||1_{[0,2^k]^d}||_q||1_{[0,2^k]^d}(x)1_{[0,4^l]}(t)|| \\
    &\leq\sum_{k,l\st k\leq 0, k+l\geq 0} 2^{-(k+l)d} 2^{kd/q}||1_{[0,2^k]^d}(x)1_{[0,4^l]}(t)||  \\
    &\leq\sum_{k,l\st k\leq 0, k+l\geq 0}2^{-(k+l)d} 2^{kd/q} 2^{l(d+2)/p}
\end{align*} and thus it is enough to obtain a bound on the sum \begin{align*}
    \sum_{k,l\st k\leq 0, k+l\geq 0}2^{-(k+l)d} 2^{kd/q} 2^{l(d+2)/p}
\end{align*} We can rewrite this sum \begin{align*}
    \sum_{k,l\st k\leq 0, k+l\geq 0}2^{-(k+l)d} 2^{kd/q} 2^{l(d+2)/p} &= \sum_{k,l\st k\leq 0, k+l\geq 0}2^{l(d+2)/p - ld - kd/q} \\
    &= \sum_{k,l\st k\leq 0, k+l\geq 0}2^{-dl - kd/p + (d+2)l/p}\\
\end{align*} where the simplification comes from the fact that $1/p + 1/q = 1$. Thus 
we can rewrite this as \begin{align*}
\sum_{k,l\st k\leq 0, k+l\geq 0}2^{-dl - kd/p + (d+2)l/p}
&\leq \sum_{k,l\st k\leq 0, k+l\geq 0}2^{-dl - (d/2(d+1)\varepsilon)((d+2)l - kd)} \\
    &\leq  \sum_{k,l\st k\leq 0, k+l\geq 0}2^{-dl - (d/2(d+1)\varepsilon)(2(d+1)l)} \\
    &\leq\sum_{k,l\st k\leq 0, k+l\geq 0}2^{-dl - (d/2(d+1)\varepsilon)((d+2)l - kd)} \\
    &= \sum_{k,l\st k\leq 0, k+l\geq 0}2^{-2(d+1)l\varepsilon} \\
\end{align*} where $\varepsilon = p/d$ for shorthand. Then we can see that this sum is finite if $2(d+1)\varepsilon > 1$.
In particular by geometric sum formulas we can bound 
\begin{align*}
    \sum_{k,l\st k\leq 0, k+l\geq 0}2^{-2(d+1)l\varepsilon} &= \sum_{k\leq 0}\frac{2^{2(d+1)k\varepsilon}}{1- 2^{2(d+1)\varepsilon}}\\
    &= \frac{1}{1-2^{2(d+1)\varepsilon}}^2\\
    &< \infty
\end{align*} provided that $p > 2(d+1)/d$ as desired.

\end{solution}
\end{document}